% =====================================================================
%  frontmatter/preface.tex —— 作者前言（はしがき，原書 p.v–vi）
%  Lead 自译
% =====================================================================
\chapter*{前言}
\addcontentsline{toc}{chapter}{前言}
\markboth{前言}{}

% [原文 p.v]
近二十年来，代数几何中最主要的发现之一，是曲面的极小模型理论可以推广到高维簇。这一推广称为极小模型纲领，或者森纲领。% [Tutor Note: 日文本作「森プログラム」，即英文的 Mori's program。之所以两个名字并用，是因为纲领的纲领性框架（锥定理、收缩、翻转、终点）由森重文在 1980 年代给出，而其高维可行性的信念来自极小模型理论本身的曲面版本。]
它原本看来唯一的目标，只是构造曲面极小模型的高维版本；然而到今天，它已发展成为一件强有力的工具，在代数几何乃至其它许多领域的问题上都有应用。

迄今为止，这一纲领只在三维完成，% [Tutor Note: 三维极小模型纲领的完成以森重文 1988 年的论文 [Mor88] 为标志；本书第 6、7 章正是这套三维理论的核心。四维以上的情形在 1998 年仍是公开问题，第 7 章 §7.6 对此有概观。]
但人们已经知道，它的许多部分在所有维数上都能起作用。

本书的目标，是向读者介绍在极小模型纲领周围发展起来的一整套思想。同时，只假定读者具备代数几何的基本知识便完成这一介绍，也是我们的目标之一。

为了这个目标，我们花了相当大的努力，尽量少作引用。许多地方简化了证明；也有一些地方，我们宁可选择易懂，而不去追求证明的完美。本来有一些定理，即使所论是局部性的，在射影簇的情形证明也要容易得多，这样的例子并不少。对它们，我们以一般形式陈述定理，却只在射影簇的情形给出证明，一般情形则只列出参考文献。极小模型纲领的应用，最终大多是关于射影簇的；对这些情形，本书给出了完整的证明。

本书的英文版 \emph{Birational Geometry of Algebraic Varieties}（Cambridge Tracts in Mathematics 134）几乎与本书同时出版。除了命题 2.71 与命题 4.11(2) 在英文版中没有之外，英文版的编号与本书相对应。% [Tutor Note: 这句话对读英文原著的读者很重要：本书（日文本）比英文版多出命题 2.71 与命题 4.11(2)，其余编号两版一致。本译注本的编号一律以日文本为准。]

% [原文 p.vi]
\section*{谢辞}
\addcontentsline{toc}{section}{谢辞}

本书之所以成为现在这个样子，很大程度上要归功于两位合作者。

H. Clemens 曾是 [CKM88] 的合著者。本书的 §1.1–1.3、§2.1、§2.2、§2.4 以及 §3.1–3.5，都是对 [CKM88] 中相应各节的改写。我们要特别感谢 H. Clemens 与出版社 Astérisque，感谢他们慨然同意我们照原样使用这些内容。% [Tutor Note: [CKM88] 即 Clemens–Kollár–Mori, \emph{Higher dimensional complex algebraic geometry}, Astérisque 166 (1988)，本书第 1 章与 §2.1–2.2、§2.4、§3.1–3.5 的材料即源于此。见书末参考文献。]

A. Corti 给我们看了他未发表的手记 “Semi-stable 3-fold flips”。那份手记把（[K+92, 第 18 章] 中所说明的）[Sho92] 的证明思想用到了半稳定翻转的情形，其中几乎不依赖二维商奇点的分类就完成了证明。对三维簇来说，这只是一个较小的区别；但我们期待，随着逐步去掉证明中所用到的低维特殊结果，它会一步步接近高维翻转的证明。Corti 的手记，正是 §7.4 中不用二维商奇点分类的半稳定翻转讨论的基础。

过去五年间，我们曾在 Park City（Utah）的地域几何研究所、京都大学数理解析研究所、Utah 大学等地，就本书所涉及的话题作过系列讲演。我们要感谢听众、同事与学生，他们提出了许多观察与改进意见，这些意见对于把本书改写成现在的样子是必要的。特别地，我们从下列各位那里得到了有益的评论：D. Abramovich、A. Bertram、J. Cheah、藤野修、川又雄二郎、H. Kley、S. Kovács、桑田孝泰、Y. Lee、R. Mayer、宫冈洋一、R. Morelli、中山昇、Th. Peternell、E. Szabó、N. Tziolas。

最后，我们要感谢 NSF（资助号 DMS-9622394）、Utah 大学、文部省国际学术研究费（共同研究计划 08044078）、文部省科学研究费基盘研究 (B)（课题号 09440010）以及稻盛财团所给予的援助。

\medskip
\noindent\hfill 1998 年 8 月

\noindent\hfill János Kollár・森重文
